\begin{proof}
Fix $x\in X$.  By \cite[Proposition~5.8 and Lemma~5.10]{GJR26}, after
shrinking to an analytic neighbourhood $W_0$ of $x$, one can write
\begin{equation}
\label{eq:ch18-zero-current-independent}
                         h=\sum_{i=1}^{s}\lambda_i h_i,
\end{equation}
where the $h_i$ are integral PL functions on $W_0$ and the real numbers
$\lambda_1,\ldots,\lambda_s$ are linearly independent over $\mathbf Q$.
Choose a compact strictly affinoid neighbourhood $A\subset W_0$ of $x$ with
$x\in\operatorname{Int}(A)$.  By
\cite[Proposition~8.3 and \S8.4]{GJR26}, choose a smooth tropicalization
\[
                         F:A\longrightarrow\mathbf R^N
\]
simultaneously tropicalizing $h_1|_A,\ldots,h_s|_A$.  Thus
$h_i|_A=\varphi_i\circ F$, with $\varphi_i$ piecewise
$(\mathbf Z,\Gamma)$-affine.  Now choose an open neighbourhood $W$ of $x$
with $W\subset\operatorname{Int}(A)$, and restrict the $h_i$ and $F$ to
$W$.  Since boundarylessness and purity are local here, we may suppose that
$W$ is boundaryless and pure of dimension $n$.  We prove that
\begin{equation}
\label{eq:ch18-zero-current-integral}
                         \dprimeddprime[h_i]=0
                         \qquad(1\leq i\leq s)
\end{equation}
on a neighbourhood of $x$.

Let $T_0=[W]$ be the integration current.  By
\cite[Example~7.5.3]{CLD25}, $T_0$ is a closed skeletal current of dimension
$n$, described in each moment map by the canonical tropical calibration.
For each $i$, put
\[
                         T_i:=\dprimeddprime(h_iT_0)
                             =\dprimeddprime[h_i].
\]
By \cite[Theorem~8.10.3]{CLD25}, $T_i$ is a skeletal current of dimension
$n-1$.

We recall the arithmetic in the construction in the proof of that theorem,
using the fixed tropicalization $F$ defined on all of $W$.  Let
$P_F^{T_0}\subset W$ be the minimal
canonically calibrated skeleton describing $T_0$ relative to $F$, as in
\cite[Propositions~7.5.7--7.5.8]{CLD25}.  In particular, $F$ is an immersion
on each cell of $P_F^{T_0}$; distinct analytic cells are not identified in
this description, even if their images overlap in $F(W)$.  Choose a common
polyhedral subdivision $\mathscr C$ of $P_F^{T_0}$ on which all the
$\varphi_i\circ F$ are affine.  A top-dimensional cell $C$ carries its
canonical calibration.  By
\cite[Propositions~14.3.1--14.3.2]{CLD25}, it is rational with respect to
the integral tangent lattice, so it has the form
\[
                         \mu_C=q_C\mu_C^{\mathrm{prim}},
                         \qquad q_C\in\mathbf Q_{>0},
\]
where $\mu_C^{\mathrm{prim}}$ is the primitive lattice volume transported
through $F|_C$.  If $D$ is a codimension-one cell and $\nu_D$ is its
primitive lattice volume, the formula in the proof of
\cite[Theorem~8.10.3]{CLD25} describes $T_i$ along $D$ by
\begin{equation}
\label{eq:ch18-zero-current-calibration}
                         \mu_{i,D}=a_{i,D}\nu_D,
                         \qquad a_{i,D}\in\mathbf Q.
\end{equation}
Indeed, the analytic degree entering the canonical tropical multiplicity is
an integer, while its comparison with the primitive tangent lattice is
rational by the cited propositions.  The derivative of the integral-affine
function $\varphi_i\circ F|_C$ in a primitive normal direction is an integer,
and every lattice transition factor is rational.  With the orientation sign
used in the boundary integral, the contribution from each cell $C$ adjacent
to $D$ is therefore rational.  Summing the finitely many contributions gives
$a_{i,D}\in\mathbf Q$.  This is the coefficient of the calibrated tropical
divisor of $\varphi_i\circ F$ on the weighted skeleton.  A further common
refinement subdivides these calibrated cycles and preserves rationality.

By linearity and \eqref{eq:ch18-zero-current-independent}, the skeletal
current $\sum_i\lambda_iT_i=\dprimeddprime[h]$ is described on the same
subdivision along $D$ by
\begin{equation}
\label{eq:ch18-zero-current-sum-calibration}
                         \left(\sum_i\lambda_i a_{i,D}\right)\nu_D.
\end{equation}
The assumed current identity says that this skeletal current is zero.  Let
$|\mathscr Q|\subset W$ denote the source-side support of the resulting
calibrated codimension-one skeleton.  Consequently the calibration in
\eqref{eq:ch18-zero-current-sum-calibration} vanishes on every
codimension-one cell.  Indeed, suppose that it were nonzero on the relative
interior of some $D$.  Choose a relatively open polyhedral set
$\Omega\Subset D^\circ$.  In the source-side skeleton, choose an open
neighbourhood $V\subset W$ of $\Omega$ such that
\[
                         V\cap|\mathscr Q|\Subset D^\circ.
\]
This is possible because the source-side subdivision is locally finite and
$\Omega$ is relatively compact in the interior of its cell.  In particular,
overlapping target images create no additional source-side incidence.
By \cite[Corollary~6.2.4]{CLD25}, after shrinking $V$ there is a nonnegative
smooth function $\chi$ compactly supported in $V$ and nonzero on $\Omega$.
Choose a constant Lagerberg $(n-1,n-1)$-form whose evaluation on $\nu_D$ is
nonzero, and change its sign if necessary.  Pairing the current with the
test form obtained by multiplying its pullback by $\chi$, the source-side
formula of \cite[Definition~7.5.1]{CLD25} gives a strictly positive value, a
contradiction.  Notice that distinct analytic cells remain distinct in this
formula even if their images under $F$ overlap.  This source-side separation
is part of the skeletal-current description and its uniqueness in
\cite[Definition~7.5.1 and Propositions~7.5.7--7.5.8]{CLD25}.  Thus
\[
                         \sum_i\lambda_i a_{i,D}=0.
\]
Since $a_{i,D}\in\mathbf Q$ and the $\lambda_i$ are $\mathbf Q$-linearly
independent, $a_{i,D}=0$ for every $i$ and $D$.  Hence the calibrated
skeletal description of each $T_i$ is zero in the fixed chart $F$ and its
polyhedral refinements.

It remains to verify the full current identity, not merely a fixed-chart
pushforward.  Let $\alpha$ be any smooth compactly supported
$(n-1,n-1)$-test form in $W$.  By the compact-support localized partition of
unity of \cite[Corollary~6.2.8]{CLD25}, together with the determination of
smooth forms by localized terms in \cite[Remark~7.2.5]{CLD25}, it is a
finite sum of terms $\rho\omega$, where $\rho$ is a smooth compactly
supported function and $\omega$ is tropicalized by a moment map $H$ on an
open set $V\subset W$ containing
$\operatorname{Supp}(\rho)$.  Since the fixed $F$ is defined on all of $W$,
the common refinement $(F|_V,H)$ is defined on $V$ and still simultaneously
tropicalizes $h_1,\ldots,h_s$.  The same computation on a common subdivision
again has rational coefficients.  Since the calibrated coefficients of
$\sum_i\lambda_iT_i=0$ vanish and the $\lambda_i$ are $\mathbf Q$-linearly
independent, the coefficient of every $T_i$ on every source-side cell of
this common refinement is zero.  The defining identity in
\cite[Definition~7.5.1]{CLD25} then gives
\[
                         \langle T_i,\rho\omega\rangle=0.
\]
Summing the localized terms proves \eqref{eq:ch18-zero-current-integral}.
Notice that no analytic current is pulled back and no target representative
is required to avoid the image of the boundary of a compact analytic
domain.  Compact support is carried by $\rho$ on the open boundaryless
space, exactly as in the definition of a skeletal current.

For each $i$, apply \cite[Theorem~18.8.4]{CLD25} with $d=1$ to the trivial
line bundle with the integral PL metric determined by $h_i$.  The separated
and boundaryless hypotheses are satisfied, and
\eqref{eq:ch18-zero-current-integral} implies that $L_{h_i}(x)$ is
numerically trivial.  The residue construction is linear, so
\[
                         L_h(x)=\sum_i\lambda_iL_{h_i}(x)
\]
is numerically trivial.  Since $x$ was arbitrary,
\cite[Proposition~7.1 and Definition~7.2]{GJR26} imply that $h$ is harmonic.

Finally, harmonicity means that $h$ and $-h$ are semipositive.  By
\Cref{thm:gr-classical-pl}, both functions are classically psh.  Hence $h$
is classically pluriharmonic.
\end{proof}
